% !PRINTED 758
\weberpage{758}

Wenn wir nun in der Gleichung (12) des vorigen Paragraphen $x=0,\alpha_1,\alpha_2,\ldots,\alpha_\mu$ setzen und die Gleichung (3) anwenden, so ergiebt sich:
\weq{4}{0 = CF(0) + F(\alpha_1) + F(\alpha_2) + \cdots + F(\alpha_\mu) + CQ(0) + e^{a_1}Q(\alpha_1) + e^{a_2}Q(\alpha_2) + \cdots + e^{a_\mu}Q(\alpha_\mu),}
und wir haben die Unmöglichkeit einer solchen Gleichung bei einer geeigneten Wahl von $f(x)$ darzuthun.

Es sei wieder $p$ eine zu unbegrenztem Wachsen bestimmte Primzahl, und
\weq{5}{f(x) = \dfrac{a^{\mu p+p-1}x^{p-1}(x-\alpha_1)^p(x-\alpha_2)^p\ldots(x-\alpha_\mu)^p}{\Pi(p-1)},}
so dass $f(x)$ eine Function mit rationalen Coëfficienten ist.

Wir bilden nun, wie im vorigen Paragraphen, durch Ordnen nach Potenzen von $x$:
\[
f(x) = \dfrac{A_{p-1}x^{p-1} + A_px^p + A_{p+1}x^{p+1}\ldots}{\Pi(p-1)},
\]
worin die $A_{p-1},A_p,\ldots$ ganze rationale Zahlen sind, und
\[
A_{p-1} = (-1)^\mu a^{\mu p+p-1}\alpha_1^p\alpha_2^p\ldots\alpha_\mu^p.
\]
Wenn wir also $p$ grösser als jede der ganzen Zahlen
\[
a, a^\mu\alpha_1\alpha_2\ldots\alpha_\mu
\]
annehmen, so ist $A_{p-1}$ durch $p$ nicht theilbar, und es wird
\[
F(0) = A_{p-1} + pA_p + p(p+1)A_{p+1} + \cdots,
\]
d.\,h. eine durch $p$ nicht theilbare ganze Zahl. Ebenso ist, wenn $p$ gross genug genommen wird, $C$ durch $p$ nicht theilbar.

Andererseits ordnen wir den Zähler von $f(x)$ nach Potenzen von $a(x-\alpha_1)$ und erhalten
\[
f(x) = \dfrac{B_pa^p(x-\alpha_1)^p + B_{p+1}a^{p+1}(x-\alpha_1)^{p+1} + \cdots}{\Pi(p-1)},
\]
worin die $B_p,B_{p+1},\ldots$ zwar nicht mehr rationale, wohl aber \gesperrt{ganze algebraische Zahlen} sind, da der Zähler von $f(x)$ eine ganzzahlige Function von $ax,a\alpha_1,\ldots,a\alpha_\mu$ ist.

Hieraus folgt nun, wie im vorigen Paragraphen:
\[
F(\alpha_1) = pB_pa^p + p(p+1)B_{p+1}a^{p+1} + \cdots
\]
Bildet man auf die gleiche Weise $F(\alpha_2),\ldots,F(\alpha_\mu)$, und beachtet, dass die Summe $B_p(\alpha_1)+B_p(\alpha_2)+\cdots+B_p(\alpha_\mu)$ eine ganze rationale Zahl ist, so folgt, dass
\[
F(\alpha_1) + F(\alpha_2) + \cdots + F(\alpha_\mu)
\]
%CONT>>
